\chapter{Discussion and Major Findings}
\label{chap:discussion-and-major-findings}

\section{Discussion of Content Findings}
\label{sec:7-1-discussion-of-content-findings}

The Level 1 analysis asked whether observable article characteristics were associated with diffusion across article–agent cases, and its most consistent result concerned title–content semantic consistency. Articles whose title and body communicated a more coherent topic were associated with greater reach and with several of the cascade-shape dimensions, and this pattern was robust to how topic was represented, to how consistency was measured, and to the dependence among observations. This finding is consistent with the expectation-management logic discussed in the literature review: the title is the first cue that shapes a reader's expectations, while the body determines whether those expectations are sustained. When the two are aligned, readers may experience the article as more coherent and relevant, which can support continued reading and resharing. When the title promises one topic but the body develops another, the mismatch may shorten attention or discourage onward sharing. The robustness checks lend some additional confidence to this reading, because the two alternative title–content distance measures, which increase as the title and body diverge, were associated with reach in the opposite direction and at a comparable magnitude. The association appears to be a property of title–content alignment.

This result also adds a particular nuance to the message-content perspective. Much of that literature locates the diffusion potential of a message in its emotional arousal or in observable surface features such as length, links, or hashtags \parencite{bergermilkman2012,lahuertaotero2018}. The semantic-consistency finding points instead to the internal coherence of a message, that is, to whether its parts agree, as a further content characteristic associated with diffusion. In a firm-side marketing setting, where articles are produced to a brief and not written spontaneously, this is a managerially tractable property: coherence between a headline and its body is something a content team can review and revise before distribution.

The two image variables require a more careful joint reading, because they are mechanically related and each coefficient is conditional on the other. The number of images was the strongest of the standardised continuous content predictors across reach and several cascade-shape dimensions, whereas the binary presence of any image was negatively associated with reach once the count was held constant. These coefficients suggest that image richness, in the sense of how many images an article carries, tracked stronger diffusion more clearly than the simple presence of an image. This is plausible in a design-oriented industry, where visual elaboration is closely tied to perceived value, so more heavily illustrated articles may read as more substantive or more appealing. The negative conditional coefficient on image presence should not be read as evidence that images depress reach. It shows that the image variables must be interpreted jointly, with the positive association carried by image quantity and not by the simple distinction between an image and no image. Article length showed a positive but comparatively small association, consistent with the idea that longer articles supply more information, though length appears to be a minor correlate relative to visual richness and title–content coherence.

Re-anchoring the topic contrasts to Home Design \& Decoration, the dominant content category, shows that topic is not inert at the content level: relative to this category, Events \& Promotions, Real Estate \& Architecture, and Customer Service \& Management articles each reached further per article–agent case, so the most common topic was not the highest-reaching one. These remain differences in average per-case reach across topics, and the substantive role of topic in this study does not rest on such main effects. The dissertation's argument is that topic matters in relation to the agent context through which an article is shared, and that relationship is examined directly at the matching level. The content-level contrasts show that topics differ in how far they travel on average, but not whether a given topic travels further when shared by an agent to whom it is suited. That conditional question is addressed by the matching account.

Finally, the explanatory power of the content-only models was modest throughout, with the reach model and the cascade-shape models each accounting for only a small share of the variation in their respective outcomes. This is the most important qualification on the content findings. Observable content characteristics were systematically associated with diffusion, but they left most of the variation in article–agent diffusion unexplained. The pattern supports the view that content alone cannot account for diffusion effectiveness in an agent-based system, in which the same article may travel very differently depending on who shares it and how well its topic fits that agent's context. The content level establishes that message characteristics are part of the story while motivating the agent-characteristic and matching levels that follow. A further nuance within the cascade-shape results is worth noting: title–content consistency was associated with reach, with cascade depth, and with whether resharing occurred at all, but not with the reshare-proportion outcome. This suggests that coherence relates more to whether an article is reshared than to the proportion of cases in which resharing happens.

\section{Discussion of Agent-Characteristic Findings}
\label{sec:7-2-discussion-of-agent-characteristic-findings}

The Level 2 analysis asked whether observable agent characteristics were associated with agent-level diffusion, and its results temper any expectation that a formal role label can summarise an agent's diffusion behaviour. Of the four role categories estimated against the Field Sales baseline, only Marketing/Operations agents differed clearly and consistently: relative to Field Sales agents, they were associated with larger, deeper, more reshared, and more structurally dispersed cascades. Design and Online Sales/Customer Service agents, by contrast, resembled the baseline on most outcomes, with only scattered and mostly small differences. Formal job category explained agent-level diffusion only partially and unevenly. This is consistent with the literature on agents and network structure, which treats the sharing actor as consequential for diffusion while recognising that a single occupational label is a coarse proxy for the audience context an agent actually represents. It also anticipates the matching argument developed below: if the static role label differentiates agents only weakly, the more informative quantity may be the relationship between an agent's role and the topic of the content being shared.

The most analytically important pattern in the agent-level results is that the same role moved in opposite directions across outcomes. Marketing/Operations agents, who were associated with stronger reach and more viral cascades, simultaneously occupied lower-centrality positions on all three reconstructed centrality measures. The agent-level outcomes cannot be collapsed into a single ranking of agent performance, because the dimension on which a role looks strongest is not the dimension on which it looks most central. This divergence is consistent with the network-caution logic and with the leakage-aware modelling design. The centrality measures are reconstructed from the same realised cascades that define reach, so they describe the shape and position of observed diffusion, not an agent's ex-ante capacity to spread content. Read in that light, the divergence is informative, not contradictory: a role associated with large, broadcast-like cascades that fan out to many recipients can register as lower in out-degree centrality than a role whose smaller cascades are more tightly interconnected. One possible reading, which would require further evidence to test directly, is that Marketing/Operations sharing reaches wide but comparatively flat audiences, while other roles diffuse through denser local clusters. What the result establishes with more confidence is methodological: reconstructed centrality is a descriptor of diffusion heterogeneity and should not be interpreted as a measure of agent capability that ought to rise and fall with reach.

Gender showed few consistent associations once role and activity were taken into account. It reached conventional significance only in isolated cases and did not show a consistent pattern across outcomes. The cautious conclusion is not that an agent's gender is irrelevant in principle. It is that, in this setting and after accounting for role and activity, gender did not carry a systematic association with diffusion outcomes. Given that gender enters as an observable demographic attribute, not as a behavioural variable, the absence of a consistent association is unsurprising and should not be over-read in either direction.

The agent's activity level, although included only as a control, was the most systematic correlate across the agent-level outcomes. More active agents tended to record higher average reach and higher centrality, and at the same time lower repeat-exposure proportions. This combination is coherent: agents who share more frequently accumulate more, and more distinct, recipients, so a smaller share of their reach comes from repeated exposure of the same contacts. Because activity is a control, not a characteristic of interest, it is not given a substantive interpretation here beyond confirming that the role and gender associations were estimated net of how much each agent shared. The robustness check that re-estimated the model without the activity control showed that the Marketing/Operations contrasts did not depend on this adjustment.

Two agent-level outcomes carry limited interpretive weight. The mean Wiener-index model accounted for essentially none of the variation and produced no significant contrasts, so it is best treated as a diagnostic. Mean reading duration is more appropriately read as a weak supplementary outcome, with low explanatory power and no significant contrasts. Following the convention adopted throughout the empirical analysis, neither is read as evidence of agent-level differences. Taken together, the agent-level findings indicate that observable agent characteristics are associated with diffusion on some dimensions but not uniformly, that the direction of association depends on which dimension is examined, and that formal role alone is a limited account of agent-level diffusion. These conclusions motivate a shift from agent attributes considered in isolation to the fit between agent context and content topic.

\section{Discussion of Matching Findings}
\label{sec:7-3-discussion-of-matching-findings}

The Level 3 analysis provides the most direct test of the dissertation's central proposition, and its results are the empirical core of the argument that, in agent-based social marketing, one size does not fit all. When article topics were better aligned with the sharing agent's professional context, the resulting agent–topic cascades were, on average, larger, deeper, more frequently reshared, and more structurally dispersed, especially under the continuous \varname{MatchScore} specification. That continuous measure carried the stronger and more precisely estimated association across these dimensions, while the threshold-based \varname{ProfessionContentMatch} measure agreed with it in sign but was weaker, reaching only marginal significance on structural virality. Because the two measures were estimated in separate specifications and yet pointed in the same direction, the matching pattern does not appear to be driven solely by the particular encoding of fit. The reach association was also the most durable result in the analysis: it persisted under exposure weighting and under the restriction to better-populated agent–topic pairs, subject to the sparse-cell caution discussed below.

This finding connects the three analytical levels into a single argument. At the content level, topics differed in average reach but content characteristics alone left most of the variation in diffusion unexplained. At the agent level, formal role differentiated agents only partially, and the same role pulled in opposite directions across outcomes. The matching result resolves part of this incompleteness by relocating the relevant quantity. Stronger diffusion was associated with the correspondence between topic and role, not with either element in isolation. This is why a job-category contrast can be largely non-significant at the agent level while a topic–role fit measure is significantly associated with diffusion at the agent–topic level: the two analyses are asking different questions. The agent-level model asks whether a role diffuses more on average across all the content it happens to share. The matching model asks whether an agent diffuses more when the content it shares fits its role. An agent role need not be uniformly stronger to perform better on the subset of topics aligned with it. This conditional, within-agent–topic relationship is what the matching measures capture.

The matching result is also where the dissertation makes contact with the content–user and content–agent fit literature. \textcite{zhang2017} found that content–user fit was associated with rebroadcasting behaviour over and above content and influence considered separately, and the influencer-congruence studies \parencite{belanche2021,breves2019} reached a parallel conclusion through the language of fit between source, message, and audience. The present finding extends that logic to a firm-side, agent-based WeChat setting in which the matching relationship is defined not between a message and a recipient's revealed preferences, but between an article's topic and the professional context of the agent through which it is shared. The homophily rationale supplies the bridge. Because social ties form disproportionately among similar people \parencite{mcpherson2001}, an agent's professional role carries information about the audience the agent is likely to reach. Topic–role alignment can serve as an observable proxy for topic–audience alignment without requiring each recipient's preferences to be measured directly. The associational framing is retained deliberately. Consistent with the caution that homophily and contagion are generically confounded in observational network data \parencite{shalizi2011}, the analysis does not claim that fit causally produces diffusion. It establishes that stronger fit was systematically associated with stronger observed diffusion on the dimensions most relevant to marketing reach.

As at the agent level, the matching measures did not move in a single direction across all outcomes: better fit was associated with lower graph-level centrality, and for the continuous \varname{MatchScore} measure also with lower agent degree centrality, mirroring the agent-level reach–centrality divergence. The same interpretive discipline applies, so a negative association between fit and centrality does not undercut the positive association between fit and reach. The practical implication is that matching should be evaluated against the reach and cascade-shape dimensions that bear on marketing diffusion. The centrality dimension should be read as a descriptor of how a well-matched cascade is shaped, not as a competing verdict on whether matching helps.

These conclusions must be qualified by the way agent–topic pairs are populated. Many pairs rest on only a handful of underlying cases, and the reach finding, though it survived both exposure weighting and the restriction to pairs supported by at least ten cases, becomes harder to test as the sample is thinned. A stricter thirty-case threshold could not be supported at all, because the few pairs that met it were too concentrated within a single topic and across only two job categories to sustain the full control specification. The matching results are robust in sign and direction, but they should not be read as precisely estimated in magnitude. Within these bounds, the Level 3 evidence supports the dissertation's central claim: the effectiveness of a marketing article in this setting was associated with the fit between article and agent, which is the empirical foundation for the personalised content-assignment implications.

\section{Managerial Implications for Personalised Content Assignment}
\label{sec:7-4-managerial-implications-for-personalised-content-assignment}

The findings translate most directly into guidance on how a firm allocates marketing content across its agents. A one-size-fits-all strategy treats agents as interchangeable broadcasting channels and pushes the same article to everyone, on the assumption that wider distribution mechanically yields wider diffusion. The Level 3 evidence offers an associational case against that assumption: in this setting, similar content topics were associated with stronger diffusion when shared through better-aligned agent contexts than when they were not. The practical recommendation is one of assignment over volume. Instead of maximising how many agents share each article, a firm may improve diffusion by directing each article towards the agents whose roles align with its topic, so that renovation- or design-oriented content reaches the audiences of design-facing agents while promotional or service-oriented content is routed to the roles better suited to it.

These assignment gains operate on top of a content foundation that the firm also controls. The Level 1 results indicated that title–content coherence, the most robust of the content associations, was associated with greater reach. Visual richness, measured through the number of images, also related positively to reach. Both are properties a content team can review before distribution. Ensuring that a headline accurately reflects the body, and that articles are sufficiently and relevantly illustrated, is a low-cost step that requires no agent-level data. This content lever is complementary, not sufficient on its own. Because the content-only models explained only a modest share of the variation in diffusion, coherent and well-presented content is best understood as a baseline condition that makes good assignment worthwhile, not as a substitute for it.

Two features of the results make this recommendation more practicable than it might first appear. First, the matching measure that performed more strongly, the continuous \varname{MatchScore}, is built from observable agent attributes and article topics, not from information that becomes available only after diffusion has occurred. A firm can compute a fit score for a prospective article–agent pairing before the article is distributed. The leakage-aware design is what secures this: because the predictors are kept distinct from the reconstructed diffusion outcomes, the matching relationship can in principle inform assignment decisions ex ante, not only describe them after the fact. Second, the threshold-based \varname{ProfessionContentMatch} measure agreed with the continuous measure in direction. This suggests that even a coarse, rule-based version of matching, in which articles are routed to a small set of topic-appropriate role categories, may capture part of the benefit without requiring a finely tuned scoring system. A firm with limited analytical capacity could begin with simple topic-to-role routing rules and refine them over time.

The multi-dimensional findings also caution against judging such a strategy on a single metric. Because reach and several cascade-shape dimensions did not move together with reconstructed centrality, a firm that evaluated matching solely through a network-position measure might wrongly conclude that it was unhelpful. The dimensions most relevant to marketing reach pointed the other way. In the matching results, well-matched content was associated with broader reach and cascade growth, while reconstructed centrality moved in the opposite direction. For practice, this means that reach and cascade-shape outcomes should be prioritised when the managerial objective is marketing diffusion, while centrality should be used diagnostically. Firms evaluating agent-based marketing programmes would benefit from tracking several diffusion dimensions in parallel.

At the same time, the agent-level results temper any expectation that assignment can be reduced to selecting a few high-performing roles. Only one role category differed clearly from the baseline, the formal role label explained agent-level diffusion only partially, and the same role that achieved stronger reach occupied a lower-centrality position. Taken with the matching evidence, this implies that the more durable lever is not identifying intrinsically superior agents but improving the fit between content and whichever agents share it. The activity-level pattern adds a modest operational qualification: concentrating distribution on a small number of highly active agents may create workload and social-capacity concerns, even if these agents historically reached more distinct recipients. A matching-based policy that spreads topic-appropriate content across a wider set of suitable agents may help balance reach against this concern.

These implications should be read within the boundaries of the present study. They rest on observational data from a single interior-design firm recorded through one SaaS platform. The matching evidence is strongest for reach, is constrained by sparsely populated agent–topic cells, and remains associational. The recommendations should be understood as a structured, empirically grounded hypothesis for how personalised content assignment might improve diffusion, and as a candidate for future experimental validation. They should not be read as a settled prescription. Within those limits, the consistent thread across all three levels is that marketing effectiveness in an agent-based social system depends less on broadcasting the same content widely than on aligning content with the agents through which it travels.
